% ============================================================================
\clearpage
\section*{v4.9.0-r39 Source-Only: Chains of Different Length under Hilbert-Shift Zoom and Mixed Readouts}
\addcontentsline{toc}{section}{v4.9.0-r39 Source-Only: Chains of Different Length under Hilbert-Shift Zoom and Mixed Readouts}
\label{sec:v490-r39-chains-zoom-mixed}
% ============================================================================

\paragraph{Purpose.}
The author states (2026-10-23) that the mixed readouts were meant for the case in which the generations form a circle after
repeated zoom, and that chains of different length might give different readings.  This layer records the pre-registered test H24
(note \texttt{llm/Vorab\_Notiz\_H24\_Ketten\_Module\_Readouts.md}, commit \texttt{8c89481}, fixed before computing).  It uses the
Hilbert shift of Definition~\ref{def:qam-hilbert-shift-plan} and the bit registry \(\beta_j=2^j\) of
Definition~\ref{def:binary-mask-registry-mixed-readouts}.  Earlier layers are not rewritten.

\noindent\textbf{Source-only status.}
This v4.9.0-r39 layer is source-only.  It does not build a PDF, does not assign a new DOI,
and does not authorize a public release.  The v4.9.0 Version DOI
\href{https://doi.org/10.5281/zenodo.23534954}{10.5281/zenodo.23534954} and the Concept DOI
\href{https://doi.org/10.5281/zenodo.19210728}{10.5281/zenodo.19210728} are unchanged by it.

\subsection*{v4.9.0-r39 model}
\label{subsec:v490-r39-model}

A chain of length \(m\) is a cyclic group \(\langle g\rangle\) with \(\operatorname{ord}g=m\).  A readout of level \(j\) (bit \(2^j\)) sees
\(\langle g^{2^j}\rangle\).  A mixed mask \(S\) with \(k=\sum_{j\in S}2^j\) is read in two ways: (\(\Sigma\)) as one step,
\(\langle g^k\rangle\), or (\(\vee\)) as the joint readout \(\langle g^{2^j}:j\in S\rangle\).  The group is
\(G=\langle\tau,\rho,T,R_{\pi/2}\rangle\cong S_3\times D_4\) of v4.9.0-r37.  The identification of zoom with the step \(2^j\) is a reading.

\subsection*{v4.9.0-r39 results}
\label{subsec:v490-r39-results}

\paragraph{Theoremlet v4.9.0-r39.1.}
\begin{enumerate}[leftmargin=2em]
\item[(a)] A chain of odd length keeps its length at every level.  For the triality, \(\tau^{2^j}=\tau\) for even \(j\) and
\(\tau^{2^j}=\tau^{-1}=\rho\tau\rho\) for odd \(j\).
\item[(b)] A chain of length \(2^a\) has length \(2^{a-j}\) at level \(j<a\) and \(1\) from level \(a\) on.  In \(G\) the visible orders are
\(\{1,2,3,4,6,12\}\) at level \(0\), \(\{1,2,3,6\}\) at level \(1\), and \(\{1,3\}\) from level \(2\) on, where every \(g^{2^j}\) lies in
\(\langle\tau\rangle\).
\item[(c)] For \(m=3\) and masks over six bits: in reading (\(\Sigma\)) the \(21\) masks with \(3\mid k\), for instance \(S=\{0,1\}\), see one
generation and the other \(42\) see three; in reading (\(\vee\)) every mask sees three.
\item[(d)] For every \(m\): (\(\Sigma\)) the visible length is \(m/\gcd(m,k)\); (\(\vee\)) it is \(m/\gcd(m,2^{\min S})\), so the odd part of \(m\)
is always preserved.
\end{enumerate}

\paragraph{Proof sketch.}
Elementary: \(2^j\) is invertible modulo odd \(m\), and \(2^j\equiv(-1)^j\pmod 3\); the subgroup of \(\mathbb Z_m\) generated by a set of
steps has order \(m/\gcd\).  All statements were checked by exact enumeration, on \(G\) as matrices, for \(m\le64\) and all \(63\) masks
(script \texttt{llm/k80\_test\_h24.py}).  \(\square\)

\paragraph{Verdict and reading.}
H24a--e are confirmed; the content is the arithmetic of cyclic groups, and the physics lies only in the reading.
\begin{enumerate}[leftmargin=2em]
\item In the model, the generations do form a circle under repeated zoom, and it is the only thing left: from level \(2\) on, only the
triality circle of length three remains visible, while the frame exchange (period \(2\)) and the duality rotation (period \(4\))
disappear.  This separates the two levels left open in v4.9.0-r38: the frames at coarse zoom, the generations at deep zoom.
\item On odd levels the circle runs backwards, and the reversed circle is the frame mirror \(\rho\); every second level reads the
generations in the mirrored frame.
\item What survives zoom is the odd part of a chain length.  Chains of length \(6\) (three generations times two frames) or \(12\) show
three at deep zoom; chains of length \(2^a\) vanish.  No generation chain of length four or more arises in \(G\).
\item Mixed readouts give different readings depending on the aggregation: as one step, a third of the masks see one generation; as
a joint readout, all see three.  Definition~\ref{def:guarded-mixed-qam-readout} leaves the aggregation \(\mu\) open.
\end{enumerate}

\paragraph{Author remark and textual reading.}
The author states (2026-10-23) that the contraction discipline of Definition~\ref{def:two-sided-admissible-contractions} was
introduced so that classical tensor calculus keeps working while mixed channels can also be contracted, that it was hardly used, and
that the paper was written from one frame.  Read against r39.1: in the selector contraction \(\chi_{mn}R^mS^n\) of
Definition~\ref{def:guarded-two-sided-mixed-qam-readout}, compatible modes are composed; for shift readouts
\(R^m=g^{2^m}\), \(S^n=g^{2^n}\) the composition is \(g^{2^m+2^n}\), which is reading (\(\Sigma\)).  The Boolean mask grammar
(\(\chi_A\vee\chi_B\)) only selects the active kinds, which is reading (\(\vee\)).  Both readings are therefore present in the text, at
different places.  Combining an even and an odd level by contraction gives \(\tau^{3}=1\): a readout composed with its frame mirror
does not see the generation circle.  The front indices of Definition~\ref{def:two-sided-structural-tensor-interface} carry the frame;
in the main body they are not varied, which is the one-frame restriction.  This is a reading, not part of the pre-registered test.

\subsection*{v4.9.0-r39 status in the three frames}
\label{subsec:v490-r39-frames}

\begin{center}
\small
\begin{tabular}{@{}>{\raggedright\arraybackslash}p{0.30\textwidth}>{\raggedright\arraybackslash}p{0.62\textwidth}@{}}
\toprule
result & M / T / P\\
\midrule
r39.1(a), (c), (d) & arithmetic of chains, frame-free\\
r39.1(b) & contains the frame maps; from level \(2\) on they are invisible, so deep-zoom statements are the same in M, T and P\\
reading (2) & odd levels read the generations in the \(\rho\)-mirrored frame\\
\bottomrule
\end{tabular}
\end{center}

\subsection*{v4.9.0-r39 register entries}
\label{subsec:v490-r39-register}

\begin{center}
\small
\begin{tabular}{@{}>{\raggedright\arraybackslash}p{0.34\textwidth}>{\raggedright\arraybackslash}p{0.58\textwidth}@{}}
\toprule
statement & status after v4.9.0-r39\\
\midrule
generation circle under zoom & the only cycle of \(G\) visible from level \(2\) on (r39.1(b)); orientation alternates with the level (r39.1(a))\\
frames versus generations & separated by zoom level in the model\\
mixed readouts & one or three generations, depending on the aggregation \(\mu\); selector contraction \(=(\Sigma)\), Boolean mask \(=(\vee)\) (reading)\\
front indices of two-sided displays & carry the frame; not varied in the main body (one-frame restriction)\\
\bottomrule
\end{tabular}
\end{center}

\subsection*{v4.9.0-r39 ledger and guard update}
\label{subsec:v490-r39-ledger-guard}

\[
R_{\rm zm}^{v490\text{-}r39}=(1,1,0)
\]
for (generation circle survives zoom; frames and generations separated by level; aggregation of mixed readouts open).

The layer spends no guard:
\[
R_{\rm glob}^{(v490\text{-}r39,+)}=(0,1,1,1,1),
\qquad
\Gamma_{\rm guard}^{v490\text{-}r39}=(1,1,1,1,1),
\qquad
\beta_4=1/256\ \text{unspent}.
\]
It does not imply exact finite \(C5\) zero closure, exact all-readout terminality,
\(C6/T2\), a global photon-clock/RH theorem, or global all-readout closure.
